\begin{table*}[t]
\centering
\caption{Matched reciprocal-clique membership comparison.}
\label{tab:clique-summary}
\scriptsize
\setlength{\tabcolsep}{1pt}
\begin{tabularx}{\textwidth}{@{}l>{\raggedleft\arraybackslash}X>{\raggedleft\arraybackslash}X>{\raggedleft\arraybackslash}X>{\raggedleft\arraybackslash}X>{\raggedleft\arraybackslash}X>{\raggedleft\arraybackslash}X>{\raggedleft\arraybackslash}X>{\raggedleft\arraybackslash}X@{}}
\toprule
Primary rule & All groups & Recip. groups & Case members & Control members & Paired diff. & Case only & Control only & Exact $p$ \\
\midrule
$k\geq4, r\geq0.50$ & 36,311 & 9,362 & 11,252 (10.96\%) & 3,925 (3.82\%) & 7.14 pp & 10,678 & 3,351 & $<0.001$ \\
\bottomrule
\end{tabularx}
\par\vspace{2pt}\begin{minipage}{0.96\linewidth}\scriptsize A candidate group is a maximal group in which every pair has a citation in at least one direction. The primary rule requires at least four members, density at least 0.75, and weighted reciprocity at least 0.50. Authors in overlapping qualifying groups are counted once per subject. The exact two-sided matched test uses only discordant pairs.
\end{minipage}
\end{table*}
